% Канонический TikZ-рисунок варианта 08. Править только здесь.
\long\def\figStereoV08{%
\begin{tikzpicture}[
    scale=0.95,
    line cap=round,
    line join=round,
    visible/.style={line width=0.9pt},
    hidden/.style={line width=0.7pt,dashed,dash pattern=on 3pt off 2.2pt}
]

% -------------------------------------------------
% Вершины пирамиды
% -------------------------------------------------

\coordinate (S) at (2.20,4.90);   % вершина
\coordinate (A) at (0.85,1.65);   % левая вершина основания
\coordinate (B) at (1.65,0.15);   % передняя вершина основания
\coordinate (C) at (4.95,1.05);   % правая вершина основания

% Точка на противоположном боковом ребре SC
% Делит его в отношении 2:5 от вершины S
\coordinate (P) at ($(S)!2/7!(C)$);

% -------------------------------------------------
% Основание
% -------------------------------------------------

\draw[visible] (A)--(B)--(C);
\draw[hidden]  (A)--(C);

% -------------------------------------------------
% Боковые рёбра
% -------------------------------------------------

\draw[visible] (S)--(A);
\draw[visible] (S)--(B);
\draw[visible] (S)--(C);

% -------------------------------------------------
% Сечение плоскостью через сторону AB
% и точку P на ребре SC
% -------------------------------------------------

\draw[hidden]  (A)--(P);
\draw[visible] (B)--(P);

\end{tikzpicture}
}
